% !TEX root = ../main.tex
\begin{ouverture}[thermo]{19}{Premier principe de la thermodynamique}{Partie 3 · L'énergie : conversions et transferts}
{Énergie interne\enspace·\enspace Travail et transfert thermique\enspace·\enspace Bilan d'énergie\enspace·\enspace Capacité thermique}
\objectif{Citer les différentes contributions microscopiques à l'énergie interne d'un système.}
\objectif{Prévoir le sens d'un transfert thermique.}
\objectif{Écrire un bilan d'énergie en distinguant la variation de l'énergie du système des transferts avec l'extérieur.}
\objectif{Exploiter la relation $\Delta U=C\,\Delta T$ pour un système incompressible.}
\objectif{Effectuer l'étude énergétique d'un système thermodynamique : chauffage, calorimétrie, changement d'état.}
\prerequis{Énergies cinétique et potentielle, travail d'une force (chapitre~15).}
\prerequis{Énergie et puissance électriques, effet Joule, rendement (classe de première).}
\prerequis{Énergie massique de changement d'état (classes de seconde et de première) ; modèle du gaz parfait
(chapitre~18).}
\end{ouverture}
\essentiel
\essSection{1}{L'énergie d'un système}
\begin{definition}[Énergie totale et énergie interne]
Un système thermodynamique est constitué d'un très grand nombre d'entités. Son énergie totale est la somme de son
énergie mécanique macroscopique $E_\mathrm{m}$ (énergie cinétique et énergie potentielle de l'ensemble du système) et de son
\textbf{énergie interne} $U$, liée à ses constituants microscopiques : $E=E_\mathrm{m}+U$.
\end{definition}
\begin{propriete}[Contributions microscopiques à l'énergie interne]
L'énergie interne regroupe :
\begin{itemize}
\item l'énergie cinétique d'agitation des entités, qui augmente avec la température ;
\item les énergies potentielles d'interaction entre entités, qui changent lors d'un changement d'état ;
\item les énergies de liaison au sein des molécules et des noyaux, mises en jeu dans les transformations chimiques et
nucléaires.
\end{itemize}
On ne sait pas calculer $U$ ; seules ses variations $\Delta U$ interviennent dans les bilans.
\end{propriete}
\essSection{2}{Le premier principe}
\begin{definition}[Deux modes de transfert d'énergie]
Un système échange de l'énergie avec l'extérieur :
\begin{itemize}
\item par \textbf{travail} $W$, transfert lié à une action macroscopique : force dont le point d'application se déplace,
courant électrique ($W=U_\mathrm{él}\,I\,\Delta t$) ;
\item par \textbf{transfert thermique} $Q$, transfert désordonné à l'échelle microscopique, dû à une différence de
température.
\end{itemize}
Par convention, $W$ et $Q$ sont \textbf{positifs} lorsque le système reçoit de l'énergie, négatifs lorsqu'il en cède.
\end{definition}
\begin{propriete}[Sens d'un transfert thermique]
Spontanément, un transfert thermique s'effectue toujours du corps le plus chaud vers le corps le plus froid. Il cesse
lorsque les températures sont égales : c'est l'\textbf{équilibre thermique}.
\end{propriete}
\begin{loi}[Premier principe de la thermodynamique]
Pour un système fermé évoluant d'un état initial à un état final, la variation de son énergie totale est égale à la
somme des énergies échangées avec l'extérieur : $\Delta E_\mathrm{m}+\Delta U=W+Q$. Pour un système macroscopiquement au
repos,
\[\Delta U=W+Q.\]
\end{loi}
\schema{19-bilan}
\begin{savoirfaire}[Effectuer un bilan d'énergie]
Définir le système, son état initial et son état final. Recenser les transferts d'énergie avec l'extérieur, travail et
transfert thermique, et préciser leur signe. Écrire le premier principe, puis exprimer chaque terme. Dans un bilan, le
membre de gauche décrit le \textbf{système} ($\Delta U$, $\Delta E_\mathrm{m}$), le membre de droite les \textbf{échanges} avec
l'extérieur ($W$, $Q$).
\end{savoirfaire}
\essSection{3}{Système incompressible et calorimétrie}
\begin{definition}[Capacité thermique]
Pour un système incompressible (solide ou liquide) qui ne change pas d'état, la variation d'énergie interne est
proportionnelle à la variation de température :
\[\Delta U=C\,\Delta T=m\,c\,\Delta T,\]
où $C$ est la capacité thermique du système (en \unit{J.K^{-1}}) et $c$ sa capacité thermique massique (en
\unit{J.kg^{-1}.K^{-1}}). L'eau a une capacité thermique massique particulièrement élevée :
$c_\mathrm{eau}=\qty{4180}{J.kg^{-1}.K^{-1}}$. Un écart de température a la même valeur en kelvins et en degrés Celsius.
\end{definition}
\begin{exemple}
Au pied d'une chute d'eau de \qty{51}{m}, comme celles du Niagara, l'énergie potentielle perdue par l'eau, une fois
l'écoulement freiné, se retrouve en énergie interne : $m\,g\,h=m\,c\,\Delta T$, d'où $\Delta T=\dfrac{g\,h}{c}=\qty{0,12}{K}$.
Joule avait imaginé mesurer cet échauffement lors de son voyage de noces.
\end{exemple}
\begin{propriete}[Changement d'état]
Lors d'un changement d'état à température constante, le système reçoit ou cède le transfert thermique
$Q=m\,L$, où $L$ est l'énergie massique de changement d'état. Pour la fusion de la glace, $L_\mathrm{f}=\qty{3,34e5}{J.kg^{-1}}$ :
faire fondre un glaçon demande autant d'énergie que chauffer la même masse d'eau liquide de \qty{80}{\degreeCelsius}.
\end{propriete}
\begin{savoirfaire}[Exploiter une mesure calorimétrique]
Dans un calorimètre, les échanges avec l'extérieur sont négligeables ($W=0$, $Q=0$) : $\Delta U_\mathrm{total}=0$. La somme des
variations d'énergie interne de tous les corps, calorimètre compris, est nulle :
\[\sum m_i\,c_i\,(T_\mathrm{f}-T_i)+C_\mathrm{cal}\,(T_\mathrm{f}-T_\mathrm{cal})=0.\]
Ainsi, \qty{200}{g} d'eau à \qty{20}{\degreeCelsius} mélangés à \qty{100}{g} d'eau à \qty{80}{\degreeCelsius} atteignent
\qty{40}{\degreeCelsius}, si l'on néglige la capacité thermique du calorimètre.
\end{savoirfaire}
\schema{19-calorimetre}
\begin{codecote}[0.64]
\begin{pythoncol}
corps = [("eau", 0.200, 4180, 20.0),
         ("eau chaude", 0.100, 4180, 80.0),
         ("cuivre", 0.150, 385, 100.0)]
C = sum(m * c for (nom, m, c, T) in corps)
Tf = sum(m*c*T for (nom, m, c, T) in corps) / C
print(round(Tf, 1))
for (nom, m, c, T) in corps:
    print(nom, round(m * c * (Tf - T)))
\end{pythoncol}
\begin{consolecol}
42.6
eau 18928
eau chaude -15616
cuivre -3312
\end{consolecol}
\cote
\fontsize{9.5}{12.5}\selectfont
\capacite{Capacité numérique}\par\smallskip
Chaque corps est décrit par son nom, sa masse (kg), sa capacité thermique massique (\unit{J.kg^{-1}.K^{-1}}) et sa
température initiale (\unit{\degreeCelsius}). La condition $\sum m_i\,c_i\,(T_\mathrm{f}-T_i)=0$ donne
$T_\mathrm{f}=\dfrac{\sum m_i\,c_i\,T_i}{\sum m_i\,c_i}$. Le programme affiche aussi l'énergie reçue par chaque corps, en
joules : leur somme est nulle, aux arrondis près.
\end{codecote}
\begin{attention}
Ne pas confondre température et transfert thermique. Un système peut recevoir un transfert thermique sans que sa
température varie (fusion de la glace), et sa température peut augmenter sans transfert thermique, par travail
(frottements, compression rapide d'un gaz).
\end{attention}
\begin{experience}[Mesurer une capacité thermique massique]
Dans un calorimètre contenant une masse d'eau connue, on plonge une résistance chauffante parcourue par un courant
$I$ sous une tension $U$ : pendant la durée $\Delta t$, l'eau et le calorimètre reçoivent le travail électrique
$W=U\,I\,\Delta t$. En mesurant l'élévation de température, on en déduit la capacité thermique de l'ensemble, puis celle de
l'eau. La méthode des mélanges permet de même de mesurer la capacité thermique massique d'un métal chauffé au
préalable dans de l'eau bouillante.
\end{experience}
\begin{rediger}
Commencer tout bilan par « Système : \dots ; état initial : \dots ; état final : \dots ». Écrire $\Delta U=W+Q$, puis
justifier le signe de chaque terme : « le système reçoit le travail électrique, $W>0$ » ; « il cède de l'énergie à l'air,
$Q<0$ ».
\end{rediger}
\exercices

\rubrique{Découverte}
\begin{blocLivre}{ex-19-1}
\begin{exercice}[1]{L'énergie interne}
\begin{questions}
\item Citer les contributions microscopiques à l'énergie interne d'un système.
\item Une balle de tennis est lancée à \qty{50}{m.s^{-1}}. Son énergie interne est-elle modifiée par ce lancer ? Quelle
forme d'énergie a augmenté ?
\item Pour chaque situation, indiquer quelle contribution à l'énergie interne varie : de l'eau qu'on chauffe de 20 à
\qty{60}{\degreeCelsius} ; un glaçon qui fond à \qty{0}{\degreeCelsius} ; du méthane qui brûle.
\item Pourquoi ne calcule-t-on jamais la valeur de l'énergie interne, mais seulement ses variations ?
\end{questions}
\end{exercice}
\end{blocLivre}
\begin{blocLivre}{ex-19-2}
\begin{exercice}[1]{Dans quel sens ?}
\begin{questions}
\item On tient dans la main une canette à \qty{5}{\degreeCelsius}. Dans quel sens s'effectue le transfert thermique ? Préciser
son signe pour le système \{canette\}, puis pour le système \{main\}.
\item Une tasse de thé à \qty{80}{\degreeCelsius} est posée dans une pièce à \qty{20}{\degreeCelsius}. Quel est le signe du
transfert thermique reçu par le thé ? Quand cesse-t-il ?
\item On se frotte les mains : elles se réchauffent. S'agit-il d'un travail ou d'un transfert thermique ?
\item Dire « le froid entre par la fenêtre » est-il correct du point de vue de la physique ? Reformuler.
\end{questions}
\end{exercice}
\end{blocLivre}
\begin{blocLivre}{ex-19-3}
\begin{exercice}[1]{Chauffer de l'eau ou de l'aluminium}
On veut porter \qty{2,0}{kg} d'eau de \qty{15}{\degreeCelsius} à \qty{60}{\degreeCelsius}.
\begin{donnees}
$c_\mathrm{eau}=\qty{4180}{J.kg^{-1}.K^{-1}}$ ; $c_\mathrm{Al}=\qty{897}{J.kg^{-1}.K^{-1}}$.
\end{donnees}
\begin{questions}
\item Calculer la variation d'énergie interne de l'eau.
\item Un chauffage de \qty{1500}{W} transmet toute son énergie à l'eau. Combien de temps faut-il ?
\item Calculer l'énergie nécessaire pour porter \qty{2,0}{kg} d'aluminium de \qty{15}{\degreeCelsius} à \qty{60}{\degreeCelsius}.
\item Pourquoi utilise-t-on de l'eau dans les bouillottes et les circuits de chauffage central ?
\end{questions}
\end{exercice}
\end{blocLivre}
\begin{blocLivre}{ex-19-4}
\begin{exercice}[1]{La bouilloire}
Une bouilloire électrique de puissance \qty{2200}{W} porte \qty{1,0}{L} d'eau de \qty{20}{\degreeCelsius} à \qty{100}{\degreeCelsius}
en \qty{3}{min}.
\begin{donnees}
$c_\mathrm{eau}=\qty{4180}{J.kg^{-1}.K^{-1}}$ ; masse volumique de l'eau : \qty{1,0}{kg.L^{-1}}.
\end{donnees}
\begin{questions}
\item Le système est l'eau. Calculer sa variation d'énergie interne.
\item Calculer le travail électrique reçu par la bouilloire, que l'on suppose entièrement transmis à l'eau.
\item Appliquer le premier principe pour calculer le transfert thermique $Q$ reçu par l'eau. Interpréter son signe.
\item Calculer le rendement du chauffage.
\end{questions}
\end{exercice}
\end{blocLivre}
\begin{blocLivre}{ex-19-5}
\begin{exercice}[1]{Faire fondre un glaçon}
Un glaçon de \qty{30}{g} à \qty{0}{\degreeCelsius} est posé sur une table.
\begin{donnees}
$L_\mathrm{f}=\qty{3,34e5}{J.kg^{-1}}$ ; $c_\mathrm{eau}=\qty{4180}{J.kg^{-1}.K^{-1}}$.
\end{donnees}
\begin{questions}
\item Quelle énergie le glaçon doit-il recevoir pour fondre entièrement ? Sa température change-t-elle pendant la
fusion ?
\item Quelle énergie faut-il ensuite fournir pour porter l'eau de fonte de \qty{0}{\degreeCelsius} à \qty{20}{\degreeCelsius} ?
\item Comparer les deux résultats.
\item L'air de la pièce et la table cèdent au glaçon une puissance moyenne de \qty{5,0}{W}. Estimer la durée de la fusion.
\end{questions}
\end{exercice}
\end{blocLivre}
\rubrique{Entraînement}
\begin{blocLivre}{ex-19-6}
\begin{exercice}[2]{Le freinage d'une voiture}
Une voiture de \qty{1200}{kg} roulant à \qty{90}{km.h^{-1}} freine jusqu'à l'arrêt sur une route horizontale. Ses quatre
disques de frein en acier ont chacun une masse de \qty{6,0}{kg}.
\begin{donnees}
$c_\mathrm{acier}=\qty{460}{J.kg^{-1}.K^{-1}}$ ; $g=\qty{9,81}{N.kg^{-1}}$.
\end{donnees}
\begin{questions}
\item Calculer l'énergie cinétique perdue par la voiture.
\item On suppose que toute cette énergie est convertie en énergie interne des disques, pendant un freinage trop bref
pour qu'ils cèdent de l'énergie à l'air. Calculer l'élévation de leur température.
\item La voiture descend ensuite un col de \qty{1000}{m} de dénivelé à vitesse constante, en utilisant seulement les freins.
Quelle énergie les freins doivent-ils dissiper ?
\item Quelle serait l'élévation de température des disques sans transfert thermique vers l'air ? Pourquoi conseille-t-on
d'utiliser le frein moteur en descente ?
\end{questions}
\end{exercice}
\end{blocLivre}
\begin{blocLivre}{ex-19-7}
\begin{exercice}[2]{Identifier un métal}
Un calorimètre de capacité thermique $C_\mathrm{cal}=\qty{40}{J.K^{-1}}$ contient \qty{200}{g} d'eau à \qty{18,0}{\degreeCelsius}. On y
plonge un bloc de métal de \qty{100}{g}, sorti d'un bain d'eau à \qty{98,0}{\degreeCelsius}. La température finale vaut
\qty{21,4}{\degreeCelsius}.
\begin{donnees}
$c_\mathrm{eau}=\qty{4180}{J.kg^{-1}.K^{-1}}$ ; capacités thermiques massiques, en \unit{J.kg^{-1}.K^{-1}} : aluminium 897 ; fer 449 ;
cuivre 385 ; plomb 129.
\end{donnees}
\begin{questions}
\item Pourquoi peut-on considérer que l'énergie interne de l'ensemble \{eau, calorimètre, métal\} ne varie pas ?
\item Écrire le bilan d'énergie et en déduire la capacité thermique massique du métal.
\item Identifier le métal et calculer l'écart relatif avec la valeur tabulée.
\item Quelle valeur aurait-on obtenue en oubliant le calorimètre ? Conclure.
\end{questions}
\end{exercice}
\end{blocLivre}
\begin{blocLivre}{ex-19-8}
\begin{exercice}[2]{Calculer une température d'équilibre en Python}
On reprend le programme de l'essentiel.
\begin{questions}
\item \capacite{Capacité numérique}Que contient la liste \code{corps} ? Que calculent les variables \code{C} et \code{Tf} ?
\item Justifier l'expression de $T_\mathrm{f}$ à partir du bilan d'énergie. Pourquoi la somme des trois énergies affichées
est-elle nulle ?
\item Le mélange est réalisé dans un calorimètre de capacité thermique \qty{40}{J.K^{-1}}, initialement à \qty{20}{\degreeCelsius}.
Quelle ligne faut-il ajouter à la liste ? Calculer la nouvelle température finale.
\item On ajoute un glaçon à \qty{0}{\degreeCelsius}. Pourquoi le programme ne convient-il plus ?
\end{questions}
\end{exercice}
\end{blocLivre}
\begin{blocLivre}{ex-19-9}
\begin{exercice}[2]{Le chauffe-eau électrique}
Un chauffe-eau électrique de \qty{200}{L} chauffe l'eau de \qty{15}{\degreeCelsius} à \qty{60}{\degreeCelsius} pendant les heures
creuses, grâce à une résistance de \qty{2200}{W}.
\begin{donnees}
$c_\mathrm{eau}=\qty{4180}{J.kg^{-1}.K^{-1}}$ ; $\qty{1}{kWh}=\qty{3,6e6}{J}$ ; prix de l'électricité en heures creuses :
0,20~€ par kilowattheure.
\end{donnees}
\begin{questions}
\item Calculer l'énergie que doit recevoir l'eau, en joules puis en kilowattheures.
\item Calculer la durée du chauffage, en supposant les pertes négligeables. Est-elle compatible avec une période de
huit heures creuses ?
\item Calculer le coût de ce chauffage.
\item Le ballon cède à la pièce \qty{1,5}{kWh} par jour lorsqu'on ne puise pas d'eau. De combien la température de l'eau
baisse-t-elle en 24~heures ?
\end{questions}
\end{exercice}
\end{blocLivre}
\begin{blocLivre}{ex-19-10}
\begin{exercice}[2]{Rafraîchir une boisson}
On veut refroidir \qty{330}{mL} de soda, assimilé à de l'eau, de \qty{25}{\degreeCelsius} à \qty{5}{\degreeCelsius}, avec des glaçons de
\qty{25}{g} à \qty{0}{\degreeCelsius}. Le verre est isolé ; on néglige sa capacité thermique.
\begin{donnees}
$c_\mathrm{eau}=\qty{4180}{J.kg^{-1}.K^{-1}}$ ; $L_\mathrm{f}=\qty{3,34e5}{J.kg^{-1}}$.
\end{donnees}
\begin{questions}
\item Quelles transformations subit la glace jusqu'à l'équilibre ?
\item Écrire le bilan d'énergie et calculer la masse de glace nécessaire. Combien de glaçons faut-il ?
\item On met trois glaçons. Calculer la température finale.
\item Avec \qty{75}{g} d'eau liquide à \qty{0}{\degreeCelsius} au lieu des glaçons, quelle serait la température finale ? Conclure.
\end{questions}
\end{exercice}
\end{blocLivre}
\begin{blocLivre}{ex-19-11}
\begin{exercice}[2]{La balle qui s'écrase}
Une balle de plomb de \qty{10}{g}, à \qty{20}{\degreeCelsius}, frappe une plaque d'acier à \qty{350}{m.s^{-1}} et s'y immobilise.
\begin{donnees}
$c_\mathrm{Pb}=\qty{129}{J.kg^{-1}.K^{-1}}$ ; température de fusion du plomb : \qty{327}{\degreeCelsius} ; énergie massique de fusion :
\qty{2,30e4}{J.kg^{-1}}.
\end{donnees}
\begin{questions}
\item Calculer l'énergie cinétique de la balle, puis son énergie cinétique par kilogramme.
\item On admet que la moitié de cette énergie se retrouve en énergie interne de la balle. Calculer sa température finale.
\item Si toute l'énergie cinétique allait à la balle, montrer que celle-ci atteindrait sa température de fusion.
\item Dans ce cas, quelle fraction du plomb fondrait ?
\end{questions}
\end{exercice}
\end{blocLivre}
\rubrique{Objectif bac}
\begin{blocLivre}{ex-19-12}
\begin{exercice}[3]{Induction ou gaz ?}
\begin{docu}{1}{Deux façons de chauffer}
On porte \qty{2,0}{L} d'eau de \qty{20}{\degreeCelsius} à \qty{100}{\degreeCelsius}. Une plaque à induction de \qty{2800}{W} y parvient
en 4~min~40~s. Un brûleur à gaz transmet à l'eau seulement 40~\% de l'énergie libérée par la combustion du méthane.
\end{docu}
\begin{docu}{2}{Données énergétiques}
Pouvoir calorifique du méthane : \qty{50}{MJ.kg^{-1}} ; combustion : $\ce{CH4 + 2 O2 -> CO2 + 2 H2O}$. En France, produire
\qty{1}{kWh} d'électricité émet en moyenne de l'ordre de \qty{50}{g} de \ce{CO2}.
\end{docu}
\begin{donnees}
$c_\mathrm{eau}=\qty{4180}{J.kg^{-1}.K^{-1}}$ ; $M(\ce{CH4})=\qty{16}{g.mol^{-1}}$ ; $M(\ce{CO2})=\qty{44}{g.mol^{-1}}$ ;
$\qty{1}{kWh}=\qty{3,6e6}{J}$.
\end{donnees}
\begin{questions}
\item Calculer l'énergie que doit recevoir l'eau.
\item Calculer le rendement de la plaque à induction.
\item \initiative Calculer la masse de méthane brûlée par le brûleur, puis la masse de dioxyde de carbone rejetée.
\item Comparer avec les émissions liées à la plaque à induction et conclure.
\end{questions}
\end{exercice}
\end{blocLivre}
\begin{blocLivre}{ex-19-13}
\begin{exercice}[3]{Stocker la chaleur du Soleil}
\begin{docu}{1}{Une centrale solaire à sels fondus}
Dans une centrale solaire thermodynamique, des miroirs concentrent le rayonnement solaire sur des tubes où circule un
fluide. Le jour, une partie de l'énergie récoltée chauffe un réservoir de sels fondus, qui la restitue la nuit pour
produire de la vapeur et faire tourner une turbine.
\end{docu}
\begin{docu}{2}{Caractéristiques, inspirées de la centrale d'Andasol (Espagne)}
Masse de sels : \num{28500} tonnes ; température des sels : \qty{291}{\degreeCelsius} (réservoir froid) et \qty{384}{\degreeCelsius}
(réservoir chaud) ; capacité thermique massique des sels : \qty{1500}{J.kg^{-1}.K^{-1}} ; rendement de la conversion en
électricité : 37,5~\% ; puissance électrique : \qty{50}{MW}, pendant 7,5~heures sans soleil selon l'exploitant.
\end{docu}
\begin{donnees}
$\qty{1}{MWh}=\qty{3,6e9}{J}$.
\end{donnees}
\begin{questions}
\item Les sels restent liquides. Justifier l'utilisation de la relation $\Delta U=m\,c\,\Delta T$.
\item Calculer l'énergie stockée lorsque tous les sels passent de \qty{291}{\degreeCelsius} à \qty{384}{\degreeCelsius}, en joules puis en
mégawattheures.
\item \initiative Pendant combien de temps la centrale peut-elle produire \qty{50}{MW} grâce à cette réserve ? Comparer à la
valeur annoncée et proposer une explication à l'écart.
\item Quel intérêt présente ce stockage pour le réseau électrique ?
\end{questions}
\end{exercice}
\end{blocLivre}
\begin{blocLivre}{ex-19-14}
\begin{exercice}[3]{Pourquoi transpire-t-on ?}
\begin{docu}{1}{Un moteur qui chauffe}
Pendant une course, les muscles d'un coureur de \qty{70}{kg} convertissent l'énergie chimique des nutriments en travail
mécanique avec un rendement d'environ 25~\%. Le reste est libéré sous forme d'énergie interne : environ \qty{900}{W} doivent
être évacués pour que la température du corps reste proche de \qty{37}{\degreeCelsius}.
\end{docu}
\begin{docu}{2}{La sueur}
L'évaporation de la sueur à la surface de la peau prélève au corps \qty{2,4e6}{J} par kilogramme évaporé. Au-delà de
\qty{39}{\degreeCelsius}, le risque de coup de chaleur devient important.
\end{docu}
\begin{donnees}
Capacité thermique massique moyenne du corps humain : \qty{3500}{J.kg^{-1}.K^{-1}}.
\end{donnees}
\begin{questions}
\item Quelle énergie le corps doit-il évacuer en une heure de course ?
\item Si cette énergie n'était pas évacuée, de combien la température du corps augmenterait-elle en une heure ?
\item \initiative Au bout de combien de minutes le coureur atteindrait-il \qty{39}{\degreeCelsius} ?
\item On suppose que toute l'énergie est évacuée par évaporation de la sueur. Quelle masse d'eau le coureur perd-il par
heure ? Pourquoi doit-il boire régulièrement ?
\end{questions}
\end{exercice}
\end{blocLivre}
\begin{blocLivre}{ex-19-15}
\begin{exercice}[3]{La canette auto-chauffante}
\begin{docu}{1}{Un café chaud sans réchaud}
Une canette auto-chauffante contient un compartiment d'oxyde de calcium \ce{CaO}. En appuyant sur le fond, on y libère de
l'eau ; la réaction $\ce{CaO(s) + H2O(l) -> Ca(OH)2(s)}$ libère \qty{65,2}{kJ} par mole d'oxyde de calcium consommé.
\end{docu}
\begin{docu}{2}{Cahier des charges}
Porter \qty{250}{mL} de café, assimilé à de l'eau, de \qty{20}{\degreeCelsius} à \qty{60}{\degreeCelsius} ; 40~\% de l'énergie libérée
est perdue (canette, air).
\end{docu}
\begin{donnees}
$c_\mathrm{eau}=\qty{4180}{J.kg^{-1}.K^{-1}}$ ; $M(\ce{CaO})=\qty{56,1}{g.mol^{-1}}$ ; $M(\ce{H2O})=\qty{18,0}{g.mol^{-1}}$.
\end{donnees}
\begin{questions}
\item Quelle contribution à l'énergie interne du système chimique diminue au cours de la réaction ?
\item Calculer l'énergie que doit recevoir le café.
\item \initiative Calculer la masse d'oxyde de calcium à prévoir.
\item Quelle masse d'eau doit réagir au minimum ? Pourquoi en met-on davantage ?
\end{questions}
\end{exercice}
\end{blocLivre}
